\chapter{Conclusions}
\label{chap:conclusions-en}

This work began with the question of whether human discussion facilitation can be described clearly enough to be implemented as a system. The literature review divided the process into observable decisions. Based on these decisions, the implemented system receives a discussion thread, analyzes its state, decides whether intervention is warranted, and generates a candidate response when appropriate.

\section{Achievement of the objectives}

The main objective was to design and implement a system based on LLMs and context-engineering techniques for analyzing and facilitating asynchronous academic discussions from different sources. It was achieved in terms of design and implementation. The system uses the discussion thread, learning objectives when supplied, and literature-derived criteria to decide whether to intervene and to produce a candidate response. The input contract allows another source to be added through an adapter, although this was not verified with a second LMS or forum.

The experiments provide evidence of technical feasibility for that objective because the facilitation sequence ran with local and external models. Since participants contribute at different times, the system can perform its analysis in the background and the execution does not require a real-time conversational response.

The experiments document the technical execution of the system. However, the evaluation did not include students, instructors, or operating courses, so it does not establish effects on learning, instructor workload, or the capacity to support more discussions.

The first specific objective concerned a staged agent system. Each stage has one responsibility and preserves an output that can be inspected. The separation allows, for example, the behavior of a role to be changed without affecting classification. In addition, interfaces for sources, model providers, and storage allow these services to be changed without modifying the nodes. The requirement concerning reasonable time and resource consumption was not verified because the full cost of each intervention was not calculated and its latency distribution was not studied.

No comparison was run against a single call to the same model, so the available evidence does not show that the multi-agent architecture produces better interventions than a direct model call. Separating the agents provides explicit responsibilities, inspectable intermediate decisions, and the ability to modify each stage independently.

The second objective was to define a literature-grounded context-engineering strategy. Facilitation roles became the starting point for the design. Studying how a human facilitator intervenes made it possible to translate pedagogical knowledge into responsibilities, criteria, and techniques. This translation provided the basis for the implementation, but it remains an engineering interpretation. Educational specialists and classroom experience are required before it can be considered a validated pedagogical strategy.

The platform-independence objective was addressed as a technical demonstration. The common contract was used with reproducible scenarios, six anonymized historical threads, and the Open edX adapter. The proof of concept implemented data adaptation and the activation mechanism. The node sequence remains fixed in the current implementation. Integration into an instructor's routine work and safe publication in a course also remain pending.

The final objective concerned system evaluation. The two sets of experimental runs produced 46 scoreable interventions, which a general-purpose LLM rated with a mean score of 3.911 out of 5. The evaluator marked 29 cases for review, including 27 with information unsupported by the discussion thread. These scores and alerts describe the recorded outputs, although they do not replace expert judgment or establish pedagogical effectiveness.

\section{Lessons from design and implementation}

The six milestones described in the introduction were completed to different extents. Problem definition and the literature review led to the process design, while implementation produced the engine and the dashboard. The common contract and Open edX made it possible to work with the discussion-thread sources, and the runs on Idril and AWS EC2 provided the technical evaluation. In addition, both environments were deployed for experimentation and demonstration.

The most difficult part was characterising model behaviour and building the system around its responses. The evaluated configurations combined differences in the models, the way structured outputs were requested, and the thresholds, so the results do not allow the differences to be attributed to a single factor. The records show execution failures and partial results. The implementation therefore had to preserve traces, limit execution times, retain partial results, and adapt the handling of model outputs.

The two environments made it possible to study different conditions. The external models had shorter mean execution times in these runs, although differences in infrastructure, provider, and load prevent a direct comparison. Idril provided an alternative based on university infrastructure and presented compatibility and structured-generation cases that did not coincide with those observed on EC2. The full cost of the interventions was not calculated, so the evidence does not establish which option would be more sustainable for a university. The data do show that the same architecture can run with local and external models.

\section{Contribution}

This work makes three technical contributions. The first is an operational decomposition of asynchronous academic discussion facilitation. Based on the roles and criteria identified in the literature, the system separates discussion-thread characterization, the decision to intervene, the selection of a role and technique, and the generation of a candidate response.

The second contribution is the implementation of that sequence through stages that can be inspected. Each stage preserves structured output and its reasoning, while validation and retry mechanisms retain partial results when a later stage fails. Discussion-thread sources, model providers, and storage connect through interfaces.

The third contribution is the traceability mechanism used to examine the system. Each results file links the discussion thread and configuration to the intermediate decisions, response, duration, and recorded errors. The dashboard and technical records provide access to this information, and the same rubric was applied to the 46 scoreable interventions.

\section{Limitations}

The evaluation imposes the following limits on the scope of the results.

\begin{itemize}
	\item The system was not evaluated with students, instructors, or operating courses. There is therefore no evidence yet that its interventions are useful in real settings, reduce instructor workload, or allow more discussions to be supported without loss of quality.
	\item Quality assurance covered two sets of runs, eight model configurations, and 18 discussion threads. This coverage exposed failures, although it does not constitute exhaustive testing across other disciplines, thread lengths, group sizes, or model behaviors.
	\item All 18 threads were written in English, and the sample was not designed to compare cultural settings or different participation norms. The system was not examined with Spanish discussions, linguistic variation within one language, or educational communities with different conventions. The results do not establish whether classifications, decisions, and responses would remain consistent in those settings.
	\item Several decisions in the sequence rely on the literature or on design criteria that were not independently tested. These include timing thresholds, retry counts, evaluator activation, and model selection by stage. No comparison varied a single factor at a time, and no systematic parameter search measured the contribution of each decision or selected its value from observed results.
	\item The state taxonomy and facilitation roles were derived from the literature and specified from an engineering perspective. They were not validated by specialists, and a single label may simplify discussions that combine several states or evolve in an unanticipated way. In addition, the classification dimensions describe the thread before an intervention; the experiments did not measure subsequent changes in those dimensions.
	\item The assessment used one general-purpose LLM as the evaluator without calibration against formal expert annotation. It was not compared with other evaluator models or repeated to examine score stability. This evaluation does not replace pedagogical and contextual review.
	\item Model-reported confidence was not calibrated against observed outcomes and is not interpreted as a reliable measure of quality.
	\item Model outputs may vary between executions. Reproducibility of classifications, decisions, and responses under the same configuration was not studied systematically.
	\item Platform independence was examined with local files and an Open edX proof of concept. It was not verified through a second integration with another LMS or forum, and the Open edX event-activation mechanism was not tested broadly enough to establish operational reliability. The publication function was tested with simulated components and does not include prior instructor review. It may also send a note generated for the instructor to the discussion thread, so it is not considered safe for automatic publication.
	\item The experiments did not include posts containing instructions directed at the model. Because student-written text forms part of the model input, the system's behavior with this type of content is unknown.
	\item OpenRouter records cost, tokens, and provider latency for individual calls, while the experiments preserve total duration for each case. The full cost of an intervention was not calculated across all calls and retries, and its latency distribution was not studied. Whether those values remain sustainable as thread volume grows is therefore unresolved.
\end{itemize}

\section{Future work}

The assessment of the 46 interventions should be repeated with several evaluator models while keeping the same rubric and cases. The comparison should examine agreement between scores and alerts, together with the stability of their justifications. This would show how strongly the results depend on the model selected as evaluator.

Model choice should be studied separately at each stage. Each experiment would change only the model used for classification, the intervention decision, role selection, or response generation. The discussion threads, instructions, context, and models used by the other stages would remain fixed. Stage completion, schema compliance, and the evaluation of the resulting decision or response could then be compared.

Changes to the context require separate experiments. Each experiment should change one element of the context while keeping the remaining configuration fixed so that its effect can be examined.

Evaluation should also include Spanish discussion threads and educational settings with different participation conventions. This requires collections of discussion threads and reviewers familiar with each setting. The analysis should examine whether the states, intervention decision, role and technique selection, and proposed response remain appropriate.

Evaluation should also be extended with real data from educational platforms and with more recent asynchronous discussions. This would make it possible to examine the system under conditions different from those represented by the discussion threads used in this work.

The recorded executions provide evidence of the technical feasibility of a modular architecture for asynchronous academic discussion facilitation. The experiments used versioned threads and local and external models. Separately, the Open edX proof of concept adapted threads to the same input contract. It remains to be evaluated whether this support improves the student experience and allows instructors to manage more discussions without reducing teaching quality.
